\section{The actual entire theta amplitude and the finite period observation}
\label{sec:ipm-entire-source}

This calculation uses the actual packet
\(h(s)=\prod_\rho(s-\rho)^{m_\rho}\), with all complete orders,
\(d_h=\deg h\ge1\), \(\Phi_h'=h\), \(\Phi_h(0)=0\),
\(g=2\xi\), \(v_h=g/h\), and
\(\upsilon_h=j_hv_h\in E_h^\times\), where
\(E_h=\mathbb C[s]/(h)\). No root outside this actual packet is
introduced. The analytic statements have these data as their domain;
they do not assert that an off-critical packet has been found.
The finite-source map is the original injective map
\(\mathcal T_hP=P(D)F_h\), with Mellin transform \(v_hP\).
Thus all maps below on \(\mathcal T_h\mathcal P_N\) are defined
unambiguously through that injective source parametrization.

\subsection{An explicit growth bound proves the global integration domain}

Write
\[
 \psi(x)=\sum_{n\ge1}e^{-\pi n^2x},\qquad
 C_\theta=(1-e^{-3\pi})^{-1},\qquad
 \beta_R=(R+1)/2.
 \tag{IPM.1}
\]
The original theta representation, also recorded in DLMF 25.5.13--14,
gives the entire identity
\[
 g(s)=1+s(s-1)\int_1^\infty
       \bigl(x^{s/2}+x^{(1-s)/2}\bigr)\psi(x)\,\frac{dx}{x}.
 \tag{IPM.2}
\]
For completeness, the identity follows initially for \(\Re s>1\)
by integrating each Gaussian:
\(\int_0^\infty\psi(x)x^{s/2}dx/x
=\pi^{-s/2}\Gamma(s/2)\zeta(s)\).
Split the integral at one and use the original Poisson identity
\(\psi(x)=x^{-1/2}\psi(1/x)+(x^{-1/2}-1)/2\).
The elementary integral on \((0,1)\) is
\(1/(s-1)-1/s=1/(s(s-1))\). Multiplication by \(s(s-1)\)
proves (IPM.2) in that half-plane. Its right side is entire, since its
integral and every coefficient derivative converge uniformly on compact
sets by Gaussian decay. Analytic continuation proves the identity
everywhere, including \(s=0,1\).

Since \(n^2\ge1+3(n-1)\), for \(x\ge1\)
one has \(\psi(x)\le C_\theta e^{-\pi x}\). Consequently
\[
 \begin{gathered}
 |g(s)|\le\mathcal X(R):=
  1+2C_\theta R(R+1)\pi^{-\beta_R}\Gamma(\beta_R,\pi)\\
 \le 1+2C_\theta R(R+1)\pi^{-\beta_R}\Gamma(\beta_R),
 \qquad |s|\le R,\ R\ge1.
 \end{gathered}
 \tag{IPM.3}
\]
Indeed both powers of \(x\) in the integral have absolute value at most
\(x^{(R+1)/2}\), and \(|s(s-1)|\le R(R+1)\).
The upper incomplete gamma is the literal integral after \(z=\pi x\).
An elementary bound sufficient here is
\(\Gamma(\beta)\le2(2\beta)^\beta\) for \(\beta\ge1\).
To prove it, write its integrand as
\((t^{\beta-1}e^{-t/2})e^{-t/2}\), bound the first factor by
its maximum \([2(\beta-1)/e]^{\beta-1}\), and integrate the second.
At \(\beta=1\) use the value one for the zeroth power. This maximum
is at most \((2\beta)^\beta\). In particular
\(\log\mathcal X(R)=O(R\log(R+2))\); the explicit function
\(\mathcal X\), rather than a substituted constant, is retained in
all tail integrals.

Let \(R_h=\max_\rho|\rho|\). For \(|s|=r\ge\max(1,2R_h)\),
\[
 |h(s)|\ge(r/2)^{d_h},\qquad
 |v_h(s)|\le 2^{d_h}r^{-d_h}\mathcal X(r).
 \tag{IPM.4}
\]
The complete-divisor hypothesis makes \(v_h\) entire, so it is bounded
on the remaining compact disc. Cauchy's integral formula on radius-one
circles gives the same \(\exp(O(r\log(r+2)))\) growth for each fixed
derivative. The formula retains both the degree and every original root
in its compact-disc and exterior bounds.

Fix \(u\ne0\), \(t\in\mathbb C\), a branch of \(\arg u\), and
the original BC rays and oriented contours
\[
 \theta_j=\frac{\arg u+\pi+2\pi j}{d_h+1},\quad
 \ell_j(r)=re^{i\theta_j},\quad
 \Gamma_j=-\ell_0+\ell_j\quad(1\le j\le d_h),\quad
 f_t=\Phi_h-ts.
 \tag{IPM.5}
\]
Their leading real exponent is
\(-r^{d_h+1}/((d_h+1)|u|)\). If
\(f_t(s)=s^{d_h+1}/(d_h+1)+\sum_{\nu=0}^{d_h}c_\nu s^\nu\),
then, for \(r\ge1\), its full real exponent is at most
\[
 -\frac{r^{d_h+1}}{(d_h+1)|u|}
       +\frac{\sum_{\nu=0}^{d_h}|c_\nu|r^\nu}{|u|}.
 \tag{IPM.6}
\]
Combining (IPM.3)--(IPM.6) proves absolute convergence of
\[
 \mathcal I_{u,t}(a)=
  \left(\int_{\Gamma_j}e^{f_t(s)/u}a(s)\,ds\right)_{j=1}^{d_h}
 \tag{IPM.7}
\]
for \(a=v_hP\), every polynomial \(P\), every derivative used below,
and every entire quotient by \(h\) constructed below. On a compact
parameter set with \(u\ne0\), the same domination proves coefficient
differentiation and the needed contour deformations inside the fixed
decaying sectors. This is a proof of the domain for the actual theta
amplitude. It makes no assertion for arbitrary undeclared rays.

